\chapter{基于一致性增强监督权重学习的单目标自适应融合定位方法}
\label{ch:single_target_fusion}

\section{引言}
\label{sec:3_intro}

在复杂城市环境中，基于低成本无人机平台的手机直连卫星非合作终端定位任务具有明显的动态性和不确定性。一方面，建筑遮挡、街道峡谷、多径反射和非视距传播会使 RSSI、AoA 以及 PDR 等不同定位分支的误差随时间和空间持续变化；另一方面，不同定位分支在误差来源和退化方式上并不相同，固定权重融合方法难以及时反映观测质量变化，容易在 LOS/NLOS 切换、局部遮挡增强或运动状态变化时出现权重失配和轨迹抖动。因此，如何根据当前环境特征和多源定位结果之间的一致性关系自适应生成融合权重，是单目标连续定位中需要解决的关键问题。

针对上述问题，本章提出一种基于一致性增强监督权重学习的单目标自适应融合定位方法。该方法保留多源定位结果在位置层面的可解释融合形式，首先由 AoA、PDR 和 RSSI 三类定位分支分别输出二维位置估计；随后利用 RMS 时延扩展、RSSI 波动特征以及多源定位结果间的一致性距离构造物理感知特征向量；在离线阶段，借助已知真实位置反解每个样本对应的最优仿射融合权重，并以此作为监督标签训练权重预测模型；在在线定位阶段，系统仅根据当前特征预测融合权重，并对三路定位结果进行自适应加权融合。与依赖人工固定权重的方法相比，该方法能够从离线数据中学习不同环境特征下的权重变化规律；与需要在线反馈更新的调权方法相比，该方法训练目标更直接，在线阶段不依赖真实位置和参数更新过程。

本章的内容组织如下：首先，建立单目标多源融合定位问题模型，给出 AoA、PDR 和 RSSI 三路定位结果的表示方式，并构建只满足权重和为 1 的仿射加权融合模型；其次，围绕物理感知特征与一致性距离构造，介绍复杂城市场景下的 RMS 时延扩展、RSSI 方差以及多源定位结果间距离等特征；随后，利用离线真实位置反解最优融合权重，并基于 XGBoost 回归模型学习从特征向量到融合权重的映射关系；最后，通过总体定位精度、典型轨迹、不同信噪比鲁棒性和消融实验验证所提方法的有效性。

\section{单目标多源融合定位问题建模}
\label{sec:3_1}

\subsection{单目标连续定位场景与任务定义}
\label{subsec:3_1_1}

本章面向复杂城市环境中手机直连卫星场景下的地面非合作终端，研究基于低成本无人机平台的非合作式无源单目标连续定位问题。与合作式定位不同，目标通常不提供同步辅助、位置信息或协议配合，无人机平台只能依赖接收到的目标相关观测信息对其位置进行估计。因此，本章关注的并非单次静态位置解算，而是复杂环境约束下的连续轨迹定位问题，重点考察定位精度、轨迹平滑性以及对环境扰动的鲁棒性。

在复杂城市环境中，高层建筑、狭窄街道和大量反射界面会引起显著的非视距传播与多径效应，使不同定位分支的有效性随时间和空间持续变化。RSSI 易受遮挡、阴影衰落和发射功率波动影响，AoA 在强反射环境中容易偏离真实来波方向，PDR 分支虽然能够提供短时连续性约束，却会在长时间运动中产生累积漂移。由此可见，三类定位分支不存在全局固定的可靠性排序，融合权重应随环境特征和分支一致性变化而自适应调整。

图~\ref{fig:single_target_scenario} 给出了本章所考虑的单目标连续定位任务场景示意。该图主要用于说明低成本无人机平台、地面目标以及复杂城市环境之间的空间关系，为后续方法框架的展开提供场景背景。

\begin{figure}[H]
    \centering
    \includegraphics[width=0.78\textwidth]{ch3/1_单目标非法终端定位场景示意图.png}
    \caption{复杂城市环境下单目标非合作终端定位场景}
    \label{fig:single_target_scenario}
\end{figure}

在上述任务场景基础上，本文构建基于一致性增强监督权重学习的自适应融合定位总体框架，如图~\ref{fig:ch3_supervised_fusion_framework} 所示。该框架保留 AoA、PDR 和 RSSI 三类定位分支的独立输出，并在位置层完成加权融合。本文将融合权重生成转化为监督回归问题：离线阶段利用真实位置反解最优权重标签，在线阶段则由 XGBoost 回归模型根据当前物理特征和一致性距离直接预测融合权重。

\begin{figure}[H]
    \centering
    \includegraphics[width=\textwidth]{ch3/2_监督权重学习融合定位总体框架图.png}
    \caption{基于一致性增强监督权重学习的单目标自适应融合定位总体框架}
    \label{fig:ch3_supervised_fusion_framework}
\end{figure}

设第 \(k\) 个定位样本对应的真实位置为
\begin{equation}
\mathbf{p}_k^{\mathrm{th}}
=
\left[
x_k^{\mathrm{th}},\,
y_k^{\mathrm{th}}
\right]^T
\in\mathbb{R}^{2},
\label{eq:3_1_true_position}
\end{equation}
其中上标 \(\mathrm{th}\) 表示该位置仅在离线训练和实验评估阶段可用。在线定位阶段，系统无法获得 \(\mathbf{p}_k^{\mathrm{th}}\)，只能根据当前观测特征和各定位分支输出完成融合估计。

\subsection{AoA、PDR 与 RSSI 多源定位结果表示}
\label{subsec:3_1_2}

为统一不同来源的定位输出，本章将各观测分支均转化为二维平面位置估计。第 \(k\) 个样本的三路定位结果分别记为
\begin{equation}
\mathbf{p}_{k,1}
=
\mathbf{p}_{k}^{\mathrm{AoA}}
\in\mathbb{R}^{2},
\quad
\mathbf{p}_{k,2}
=
\mathbf{p}_{k}^{\mathrm{PDR}}
\in\mathbb{R}^{2},
\quad
\mathbf{p}_{k,3}
=
\mathbf{p}_{k}^{\mathrm{RSSI}}
\in\mathbb{R}^{2},
\label{eq:3_2_multi_source_position}
\end{equation}
其中 \(\mathrm{PDR}\) 表示 PDR 分支，用于提供短时运动连续性约束。若实验数据中该分支由 PDR 输出获得，则本文将其作为运动先验或 PDR 结果使用，而不要求非合作目标在在线定位阶段主动提供惯性数据。

对于 AoA 分支，目标相对于第 \(m\) 个观测位置的到达方位角可表示为
\begin{equation}
\theta_{m,k}^{\mathrm{AoA}}
=
\mathrm{atan2}(y_k-y_{m,k},\,x_k-x_{m,k})
+
b_{m,k}^{\theta}
+
v_{m,k}^{\theta},
\label{eq:3_1_aoa_branch}
\end{equation}
其中，\((x_{m,k},y_{m,k})\) 为第 \(m\) 个观测位置的平面坐标，\(b_{m,k}^{\theta}\) 表示由 NLOS 偏移、阵列失配或反射路径主导带来的角度偏差，\(v_{m,k}^{\theta}\) 为角度测量噪声。通过多观测位置的方位线交汇或等价几何求解，可得到 AoA 分支的位置估计 \(\mathbf{p}_{k}^{\mathrm{AoA}}\)。

对于 PDR 分支，本文利用短时运动连续性给出当前位置预测。若记上一时刻递推位置为 \(\mathbf{p}_{k-1}^{\mathrm{PDR}}\)，步长或位移幅度估计为 \(\hat{l}_k\)，航向角估计为 \(\hat{\psi}_k\)，则递推关系可写为
\begin{equation}
 \mathbf{p}_{k}^{\mathrm{PDR}}
=
\mathbf{p}_{k-1}^{\mathrm{PDR}}
+
\left[
\hat{l}_k\cos \hat{\psi}_k,\,
\hat{l}_k\sin \hat{\psi}_k
\right]^T
+
\mathbf{w}_k,
\label{eq:3_1_pdr_branch}
\end{equation}
其中，\(\mathbf{w}_k\in\mathbb{R}^{2}\) 为递推噪声和累计漂移误差。该分支不直接依赖当前无线传播观测，能够在无线观测短时退化时提供轨迹连续性约束，但其长期可靠性会受到累计漂移影响。

对于 RSSI 分支，接收信号强度可由带有环境附加偏差的对数距离路径损耗模型描述：
\begin{equation}
r_{m,k}
=
P_0
-
10\eta\log_{10}\!\left(\frac{d_{m,k}}{d_0}\right)
+
b_{m,k}^{\mathrm{NLOS}}
+
v_{m,k}^{r},
\label{eq:3_1_rssi_branch}
\end{equation}
其中，\(P_0\) 表示参考距离 \(d_0\) 处的接收功率，\(\eta\) 为路径损耗指数，\(d_{m,k}\) 为目标与观测位置之间的几何距离，\(b_{m,k}^{\mathrm{NLOS}}\) 表示遮挡和多径引起的附加偏差，\(v_{m,k}^{r}\) 为测量噪声。基于多个观测位置的 RSSI 测量，可通过路径损耗反演与几何约束获得 RSSI 分支的位置估计 \(\mathbf{p}_{k}^{\mathrm{RSSI}}\)。

\subsection{仿射约束下的自适应加权融合模型}
\label{subsec:3_1_3}

将式~\eqref{eq:3_2_multi_source_position} 中的三路定位结果按列拼接，可得到第 \(k\) 个样本的多源定位矩阵
\begin{equation}
\mathbf{P}_k
=
\left[
\mathbf{p}_{k,1},\,
\mathbf{p}_{k,2},\,
\mathbf{p}_{k,3}
\right]
=
\begin{bmatrix}
x_{k,1} & x_{k,2} & x_{k,3}\\
y_{k,1} & y_{k,2} & y_{k,3}
\end{bmatrix}
\in\mathbb{R}^{2\times3}.
\label{eq:3_3_position_matrix}
\end{equation}
对应的融合权重定义为
\begin{equation}
\mathbf{q}_k
=
\left[
q_{k,1},\,
q_{k,2},\,
q_{k,3}
\right]^T
\in\mathbb{R}^{3}.
\label{eq:3_4_weight_vector}
\end{equation}
因此，融合定位结果可写为
\begin{equation}
\hat{\mathbf{p}}_k^{\mathrm{Fused}}
=
\mathbf{P}_k\mathbf{q}_k
=
q_{k,1}\mathbf{p}_{k,1}
+
q_{k,2}\mathbf{p}_{k,2}
+
q_{k,3}\mathbf{p}_{k,3}
\in\mathbb{R}^{2}.
\label{eq:3_5_affine_fusion}
\end{equation}

本文对融合权重施加仿射约束
\begin{equation}
\mathbf{e}^T\mathbf{q}_k=1,
\quad
\mathbf{e}=[1,1,1]^T\in\mathbb{R}^{3},
\label{eq:3_6_affine_constraint}
\end{equation}
但不强制要求 \(q_{k,i}\ge0\)。与非负凸组合相比，仿射融合允许融合结果在三路定位结果构成的三角形外进行有限外推，从而为补偿系统性偏差提供更大的自由度。例如，当三路定位分支在某类环境中存在同向偏移时，非负加权平均难以越过三个分支结果的包络范围，而仿射权重可以通过负权重或大于 1 的权重对偏差进行修正。需要指出的是，权重自由度增加也可能带来数值不稳定问题，因此后续最优权重反解过程需要引入正则化约束，以抑制极端权重。

由式~\eqref{eq:3_5_affine_fusion} 可知，本文方法的核心不在于改变位置层融合形式，而在于根据当前物理环境和多源定位结果之间的相互关系生成合适的 \(\mathbf{q}_k\)。离线阶段利用真实位置构造最优权重标签，在线阶段则由监督学习模型根据特征向量预测权重，从而实现不依赖在线真值的自适应融合定位。

\section{物理感知特征与一致性距离构造}
\label{sec:3_3}

\subsection{复杂城市环境下的物理特征提取}
\label{subsec:3_3_1}

为了建立观测波动与环境成因之间的对应关系，本文引入基于射线追踪的高保真场景建模方法，用于构造复杂城市环境中的传播先验信息。对于本章所研究的单目标连续定位任务而言，监督权重学习模型的训练和评估需要大量带有真实位置标注的样本，同时还需要覆盖不同轨迹、不同信噪比以及不同视距/非视距切换条件。真实城市环境下的大规模采集代价较高，且难以同时保证多轨迹可控、信道状态可重复和位置真值准确。因此，本文采用“真实地理数据建模 + 射线追踪传播仿真”的方式构建训练与测试场景，以兼顾物理一致性、样本规模和实验可控性。

在场景建模阶段，本文首先利用 OpenStreetMap（OSM）导出的真实地理信息数据，获取目标区域内建筑物的平面轮廓、高度信息及空间分布关系，并在 Site Viewer 环境中完成三维城市几何场景重建。为了尽可能逼近真实传播条件，建模区域选取北京市典型建筑密集街区作为原型场景，通过将建筑物几何轮廓映射到三维坐标系中，形成包含高层建筑、街道走廊和开阔区域的城市峡谷环境。该建模过程为后续电磁传播仿真提供了几何基础，也使得不同位置的传播差异能够与真实城市结构建立对应关系。图~\ref{fig:ch3_osm_city_model} 展示了基于 OSM 数据构建的复杂城市三维场景，该图直观反映了后续定位实验所依赖的环境基础。

\begin{figure}[H]
    \centering
    \includegraphics[width=0.82\textwidth]{ch3/3_复杂城市三维场景建模示意图.jpg}
    \caption{基于 OSM 数据的复杂城市三维场景建模}
    \label{fig:ch3_osm_city_model}
\end{figure}

从图~\ref{fig:ch3_osm_city_model} 可以看出，目标区域内建筑分布具有明显的不规则性，高层建筑、道路开口和街区边界共同决定了传播条件的空间差异。这种真实几何结构使得 LOS 与 NLOS 状态切换、多径强弱变化以及阴影区分布具有明确的空间依托，从而保证后续实验中环境先验提取与定位性能分析具有较强的工程可信度。

在物理特征方面，本文首先选取 RMS 时延扩展 \(\sigma_{\tau,k}\) 表征第 \(k\) 个样本处的多径复杂度。该特征反映接收信号在时延域上的扩展程度，可由仿真平台中的路径级传播参数获得；在实际部署中，则可由宽带信道估计、同步相关峰或预先构建的环境传播先验库近似得到。若采用传播先验库，在线阶段只能依据当前初始定位结果、上一时刻融合位置或可观测链路状态进行查询，不能以真实位置 \(\mathbf{p}_k^{\mathrm{th}}\) 作为索引。需要强调的是，\(\sigma_{\tau,k}\) 不使用真实位置标签，其作用是为权重预测模型提供当前传播环境复杂度的物理描述。

设射线追踪输出 \(N_k\) 条有效传播路径，第 \(n\) 条路径的复幅度为 \(\alpha_{n,k}\)、时延为 \(\tau_{n,k}\)，则功率加权平均时延为
\begin{equation}
\bar{\tau}_k
=
\frac{\sum_{n=1}^{N_k}|\alpha_{n,k}|^2\tau_{n,k}}
{\sum_{n=1}^{N_k}|\alpha_{n,k}|^2},
\label{eq:3_7_mean_delay}
\end{equation}
RMS 时延扩展定义为
\begin{equation}
\sigma_{\tau,k}
=
\sqrt{
\frac{\sum_{n=1}^{N_k} |\alpha_{n,k}|^2(\tau_{n,k}-\bar{\tau}_k)^2}
{\sum_{n=1}^{N_k} |\alpha_{n,k}|^2}
}.
\label{eq:3_8_rms_delay}
\end{equation}
\(\sigma_{\tau,k}\) 越大，通常表示多径分量在时延域上的离散程度越强，当前无线定位结果更容易受到反射路径和非视距传播影响。

其次，为刻画接收功率的局部波动程度，本文引入 RSSI 方差 \(\sigma^2_{\mathrm{RSSI},k}\)。设滑动窗口长度为 \(W\)，窗口内 RSSI 序列为 \(\{r_i\}_{i=k_W}^{k}\)，其中 \(k_W=\max(1,k-W+1)\)，窗口均值为
\begin{equation}
\bar{r}_k
=
\frac{1}{k-k_W+1}
\sum_{i=k_W}^{k}r_i,
\label{eq:3_9_rssi_mean}
\end{equation}
则 RSSI 方差可写为
\begin{equation}
\sigma^2_{\mathrm{RSSI},k}
=
\frac{1}{k-k_W+1}
\sum_{i=k_W}^{k}
\left(
r_i-\bar{r}_k
\right)^2.
\label{eq:3_10_rssi_var}
\end{equation}
该特征反映局部时间范围内功率观测的稳定性。当遮挡、深衰落或局部散射增强时，RSSI 序列通常会出现更明显的波动，对应的 \(\sigma^2_{\mathrm{RSSI},k}\) 也会增大。与单点 RSSI 相比，窗口方差能够更稳定地刻画功率观测的局部可靠性，避免由单次随机噪声引起的误判。

\subsection{多源定位结果间的一致性距离定义}
\label{subsec:3_3_2}

仅依赖 RMS 时延扩展和 RSSI 方差等物理特征，仍然难以直接判断三路定位分支之间是否相互支持。为此，本文进一步引入多源定位结果间的一致性距离，用于刻画不同定位分支输出结果之间的空间冲突程度。对于第 \(k\) 个样本，三路定位结果分别为 \(\mathbf{p}_{k,1},\mathbf{p}_{k,2},\mathbf{p}_{k,3}\in\mathbb{R}^{2}\)，则两两一致性距离定义为
\begin{equation}
d_{12,k}
=
\left\|
\mathbf{p}_{k,1}-\mathbf{p}_{k,2}
\right\|_2,
\label{eq:3_10_consistency_12}
\end{equation}
\begin{equation}
d_{13,k}
=
\left\|
\mathbf{p}_{k,1}-\mathbf{p}_{k,3}
\right\|_2,
\label{eq:3_11_consistency_13}
\end{equation}
\begin{equation}
d_{23,k}
=
\left\|
\mathbf{p}_{k,2}-\mathbf{p}_{k,3}
\right\|_2.
\label{eq:3_12_consistency_23}
\end{equation}
其中，\(d_{12,k}\)、\(d_{13,k}\) 和 \(d_{23,k}\) 均为标量。若三路定位结果彼此接近，则说明当前多源输出具有较好一致性，融合权重可相对均衡；若某一路结果与另外两路明显分离，则说明该路定位分支可能受到局部传播退化、几何失配或累计漂移影响，需要在权重预测中降低其贡献。

一致性距离并不直接使用真实位置，因此可以同时用于离线训练和在线推理。它的作用在于将多源定位结果之间是否冲突这一结构信息显式输入权重预测模型，使模型不只根据单一路径的物理特征判断可靠性，还能根据三路输出之间的相对关系学习权重分配规律。例如，当 AoA 和 RSSI 两路无线定位结果相互接近而 PDR 分支明显偏离时，较大的 \(d_{12,k}\) 或 \(d_{23,k}\) 可以提示模型 PDR 结果可能存在累积漂移；当 RSSI 分支与其他两路均明显偏离时，则可能表明当前功率观测受到遮挡或局部衰落影响。需要指出的是，一致性距离本身只能描述分支之间的相对关系，并不能单独判断哪一路结果一定正确；因此，本文将其与 RMS 时延扩展、RSSI 波动等物理特征联合输入模型，由监督学习过程从离线样本中学习不同特征组合下的权重分配规律。

\subsection{面向权重预测的特征向量构造}
\label{subsec:3_3_3}

综合上述物理特征和一致性距离，本文构造第 \(k\) 个样本的权重预测特征向量。若采用 RMS 时延扩展、RSSI 方差、窗口平均 RSSI 以及三组一致性距离作为输入，则特征向量可写为
\begin{equation}
\mathbf{t}_k
=
\left[
\sigma_{\tau,k},\,
\sigma^2_{\mathrm{RSSI},k},\,
\bar{r}_k,\,
d_{12,k},\,
d_{13,k},\,
d_{23,k}
\right]^T,
\label{eq:3_13_feature_vector}
\end{equation}
若仅使用式~\eqref{eq:3_13_feature_vector} 中 6 个分量，则 \(\mathbf{t}_k\in\mathbb{R}^{6}\)。在实际实现中，若进一步加入 SNR、RSSI 均值、速度估计、上一时刻预测权重或其他窗口统计量，则可统一记为
\begin{equation}
\mathbf{t}_k\in\mathbb{R}^{d},
\label{eq:3_14_general_feature_dim}
\end{equation}
其中 \(d\) 为实际采用的特征维度。后续监督权重学习模型的输入即为 \(\mathbf{t}_k\)，输出为三维融合权重 \(\hat{\mathbf{q}}_k\in\mathbb{R}^{3}\)。

为保证不同量纲特征在数据记录和模型训练中的数值范围相对一致，本文在构造特征矩阵前对各维特征进行标准化处理。设第 \(j\) 维特征在训练集上的均值和标准差分别为 \(\mu_j^{\mathrm{std}}\) 和 \(\sigma_j^{\mathrm{std}}\)，则标准化后的特征为
\begin{equation}
\tilde{t}_{k,j}
=
\frac{t_{k,j}-\mu_j^{\mathrm{std}}}{\sigma_j^{\mathrm{std}}+\epsilon_{\mathrm{std}}},
\label{eq:3_15_feature_standardization}
\end{equation}
其中，\(\epsilon_{\mathrm{std}}\) 为防止分母为零的小常数。需要说明的是，树模型本身对单调尺度变换并不敏感，标准化并不是 XGBoost 训练的必要条件；本文保留该处理主要是为了统一特征记录方式，并便于与其他回归模型或消融实验保持一致。在线推理阶段使用的 \(\mu_j^{\mathrm{std}}\) 和 \(\sigma_j^{\mathrm{std}}\) 均由训练集统计得到，不使用测试样本真值信息。

\section{离线最优融合权重反解}
\label{sec:3_4}

上一节构造了用于权重预测的物理感知特征向量 \(\mathbf{t}_k\)，但监督学习模型仍需要明确的训练标签。本文的基本思路是：在离线阶段，真实位置 \(\mathbf{p}_k^{\mathrm{th}}\) 已知，因此可以针对每个样本直接求解使融合定位误差最小的权重 \(\mathbf{q}_k^\ast\)，再将 \((\mathbf{t}_k,\mathbf{q}_k^\ast)\) 作为后续权重回归模型的监督训练样本。在线阶段则不再使用真实位置，而是由训练完成的模型根据当前特征预测权重。

\subsection{基于真实位置的最优权重优化模型}
\label{subsec:3_4_1}

对于第 \(k\) 个离线样本，三路定位结果组成的矩阵为 \(\mathbf{P}_k\in\mathbb{R}^{2\times3}\)，真实位置为 \(\mathbf{p}_k^{\mathrm{th}}\in\mathbb{R}^{2}\)，融合权重为 \(\mathbf{q}_k\in\mathbb{R}^{3}\)。在给定 \(\mathbf{P}_k\) 和 \(\mathbf{p}_k^{\mathrm{th}}\) 的条件下，最直接的权重反解目标是最小化融合结果与真实位置之间的平方误差：
\begin{equation}
\min_{\mathbf{q}_k\in\mathbb{R}^{3}}
\left\|
\mathbf{P}_k\mathbf{q}_k-\mathbf{p}_k^{\mathrm{th}}
\right\|_2^2,
\quad
\mathrm{s.t.}\quad
\mathbf{e}^T\mathbf{q}_k=1.
\label{eq:3_16_unregularized_weight_inverse}
\end{equation}
式~\eqref{eq:3_16_unregularized_weight_inverse} 只施加仿射约束，不施加非负约束。其含义是，融合结果不必局限于三路定位结果构成的凸包内部，而可以通过仿射外推补偿多源定位结果共同存在的系统性偏差。若三路定位点在二维平面内位置分布良好，则式~\eqref{eq:3_16_unregularized_weight_inverse} 能够给出使当前样本融合误差最小的仿射权重。

不过，在实际数据中，三路定位结果可能接近共线，或者多个分支在某些时刻高度接近。此时，单纯求解式~\eqref{eq:3_16_unregularized_weight_inverse} 可能得到幅度过大的权重。过大的正负权重虽然能在离线样本上降低误差，但不利于后续监督模型学习稳定的权重规律。为此，本文在反解过程中引入以均匀权重为中心的二范数正则项，构造如下优化问题：
\begin{equation}
\mathbf{q}_k^\ast
=
\arg\min_{\mathbf{q}_k\in\mathbb{R}^{3}}
\left\|
\mathbf{P}_k\mathbf{q}_k-\mathbf{p}_k^{\mathrm{th}}
\right\|_2^2
+
\mu
\left\|
\mathbf{q}_k-\mathbf{q}_0
\right\|_2^2,
\quad
\mathrm{s.t.}\quad
\mathbf{e}^T\mathbf{q}_k=1,
\label{eq:3_17_regularized_weight_inverse}
\end{equation}
其中，\(\mu\ge0\) 为正则化系数，
\begin{equation}
\mathbf{q}_0
=
\frac{1}{3}\mathbf{e}
=
\left[
\frac{1}{3},\,
\frac{1}{3},\,
\frac{1}{3}
\right]^T
\in\mathbb{R}^{3}
\label{eq:3_18_uniform_weight}
\end{equation}
为均匀融合权重。正则项并不改变权重和为 1 的仿射约束，也不强制权重非负，其作用只是抑制由几何退化造成的极端权重。特别地，当 \(\mu=0\) 时，式~\eqref{eq:3_17_regularized_weight_inverse} 退化为无正则的仿射最小二乘反解。

\subsection{仿射约束最小二乘问题的 KKT 求解}
\label{subsec:3_4_2}

式~\eqref{eq:3_17_regularized_weight_inverse} 是带线性等式约束的凸二次优化问题。由于变量仅为三维权重向量 \(\mathbf{q}_k\)，可以通过 KKT 条件稳定求解。首先将目标函数写成标准二次型。令
\begin{equation}
\mathbf{A}_k
=
\mathbf{P}_k^T\mathbf{P}_k
+
\mu\mathbf{I}_3
\in\mathbb{R}^{3\times3},
\label{eq:3_19_A_matrix}
\end{equation}
\begin{equation}
\mathbf{b}_k
=
\mathbf{P}_k^T\mathbf{p}_k^{\mathrm{th}}
+
\mu\mathbf{q}_0
\in\mathbb{R}^{3}.
\label{eq:3_20_b_vector}
\end{equation}
忽略与 \(\mathbf{q}_k\) 无关的常数项，式~\eqref{eq:3_17_regularized_weight_inverse} 的目标函数可写为
\begin{equation}
J_k(\mathbf{q}_k)
=
\mathbf{q}_k^T\mathbf{A}_k\mathbf{q}_k
-
2\mathbf{b}_k^T\mathbf{q}_k
+ C_k,
\label{eq:3_21_quadratic_objective}
\end{equation}
其中，\(C_k\) 为常数项，不影响最优解。构造拉格朗日函数
\begin{equation}
\mathcal{L}(\mathbf{q}_k,\lambda_k)
=
\mathbf{q}_k^T\mathbf{A}_k\mathbf{q}_k
-
2\mathbf{b}_k^T\mathbf{q}_k
+
\lambda_k
\left(
\mathbf{e}^T\mathbf{q}_k-1
\right),
\label{eq:3_22_lagrangian}
\end{equation}
其中，\(\lambda_k\in\mathbb{R}\) 为等式约束对应的拉格朗日乘子。令 \(\mathcal{L}\) 对 \(\mathbf{q}_k\) 和 \(\lambda_k\) 的偏导为零，可得 KKT 方程
\begin{equation}
\begin{bmatrix}
2\mathbf{A}_k & \mathbf{e}\\
\mathbf{e}^T & 0
\end{bmatrix}
\begin{bmatrix}
\mathbf{q}_k^\ast\\
\lambda_k
\end{bmatrix}
=
\begin{bmatrix}
2\mathbf{b}_k\\
1
\end{bmatrix}.
\label{eq:3_23_kkt_system}
\end{equation}
在维度上，
\begin{equation}
\begin{bmatrix}
2\mathbf{A}_k & \mathbf{e}\\
\mathbf{e}^T & 0
\end{bmatrix}
\in\mathbb{R}^{4\times4},
\quad
\begin{bmatrix}
\mathbf{q}_k^\ast\\
\lambda_k
\end{bmatrix}
\in\mathbb{R}^{4},
\quad
\begin{bmatrix}
2\mathbf{b}_k\\
1
\end{bmatrix}
\in\mathbb{R}^{4}.
\label{eq:3_24_kkt_dimension}
\end{equation}
当 \(\mu>0\) 时，\(\mathbf{A}_k\) 为正定矩阵，且约束向量 \(\mathbf{e}\neq\mathbf{0}\)，上述 KKT 系统具有唯一解。当 \(\mu=0\) 时，若三路定位点在仿射意义下退化，例如三点接近共线或局部高度重合，则无正则反解可能出现不唯一或数值不稳定。因此，实际实现中本文采用 \(\mu>0\) 的正则化形式，并且不显式求矩阵逆，而是直接求解式~\eqref{eq:3_23_kkt_system} 所示的线性方程组，以降低病态样本对数值稳定性的影响。解得的前三个分量即为
\begin{equation}
\mathbf{q}_k^\ast
=
\left[
q_{k,1}^\ast,\,
q_{k,2}^\ast,\,
q_{k,3}^\ast
\right]^T
\in\mathbb{R}^{3},
\label{eq:3_25_optimal_weight_label}
\end{equation}
并满足
\begin{equation}
\mathbf{e}^T\mathbf{q}_k^\ast=1.
\label{eq:3_26_optimal_weight_constraint}
\end{equation}

\subsection{权重标签生成与反解质量检查}
\label{subsec:3_4_3}

对所有离线样本重复上述反解过程，即可得到每个样本对应的最优权重标签。设离线数据集中共有 \(N\) 个样本，则第 \(k\) 个样本的标签生成过程可概括为：首先由三路定位结果构造 \(\mathbf{P}_k\in\mathbb{R}^{2\times3}\)，并由真实位置 \(\mathbf{p}_k^{\mathrm{th}}\in\mathbb{R}^{2}\) 构造式~\eqref{eq:3_17_regularized_weight_inverse}；随后求解式~\eqref{eq:3_23_kkt_system}，得到 \(\mathbf{q}_k^\ast\in\mathbb{R}^{3}\)；最后将其与同一样本的特征向量配对，形成监督学习样本。

根据第~\ref{subsec:3_3_3} 小节的特征定义，本文在实验中采用 6 维输入特征，即
\begin{equation}
\mathbf{t}_k
=
\left[
\sigma_{\tau,k},\,
\sigma^2_{\mathrm{RSSI},k},\,
\bar{r}_k,\,
d_{12,k},\,
d_{13,k},\,
d_{23,k}
\right]^T
\in\mathbb{R}^{6}.
\label{eq:3_32_final_feature_vector}
\end{equation}
其中前三个分量反映传播环境和接收功率波动，后三个分量反映 AoA、PDR 和 RSSI 三路定位结果之间的一致性关系。需要注意的是，式~\eqref{eq:3_32_final_feature_vector} 中的所有分量均可由当前观测、历史窗口统计或三路定位结果计算得到，不包含真实位置 \(\mathbf{p}_k^{\mathrm{th}}\)。真实位置只用于式~\eqref{eq:3_17_regularized_weight_inverse} 中生成监督标签。

对 \(N\) 个样本按行堆叠，可得到特征矩阵
\begin{equation}
\mathbf{T}
=
\begin{bmatrix}
\mathbf{t}_1^T\\
\mathbf{t}_2^T\\
\vdots\\
\mathbf{t}_N^T
\end{bmatrix}
\in\mathbb{R}^{N\times 6},
\label{eq:3_27_feature_matrix}
\end{equation}
以及权重标签矩阵
\begin{equation}
\mathbf{Q}^{\ast}
=
\begin{bmatrix}
\left(\mathbf{q}_1^\ast\right)^T\\
\left(\mathbf{q}_2^\ast\right)^T\\
\vdots\\
\left(\mathbf{q}_N^\ast\right)^T
\end{bmatrix}
\in\mathbb{R}^{N\times3}.
\label{eq:3_28_label_matrix}
\end{equation}
其中，\(\mathbf{T}\) 的每一行对应一个样本的 6 维特征，\(\mathbf{Q}^{\ast}\) 的每一行对应同一样本的三维最优仿射权重。二者共同构成后续监督权重学习模型的训练数据：
\begin{equation}
\mathcal{D}_{\mathrm{w}}
=
\left\{
\left(\mathbf{t}_k,\mathbf{q}_k^\ast\right)
\right\}_{k=1}^{N}.
\label{eq:3_29_weight_learning_dataset}
\end{equation}
从矩阵维度看，本文后续回归模型的输入为 \(\mathbf{T}\in\mathbb{R}^{N\times6}\)，输出标签为 \(\mathbf{Q}^{\ast}\in\mathbb{R}^{N\times3}\)，因此权重学习任务可以表示为从 6 维特征空间到 3 维连续权重空间的映射学习问题：
\begin{equation}
f:\mathbb{R}^{6}\rightarrow\mathbb{R}^{3},
\qquad
\hat{\mathbf{q}}_k=f(\mathbf{t}_k).
\label{eq:3_33_weight_mapping}
\end{equation}

为了检查权重反解结果是否稳定，本文进一步记录反解后的最小融合误差
\begin{equation}
e_k^\ast
=
\left\|
\mathbf{P}_k\mathbf{q}_k^\ast
-
\mathbf{p}_k^{\mathrm{th}}
\right\|_2,
\label{eq:3_30_inverse_error}
\end{equation}
以及权重幅度
\begin{equation}
R_{q,k}
=
\left\|
\mathbf{q}_k^\ast
\right\|_2.
\label{eq:3_31_inverse_weight_norm}
\end{equation}
其中，\(e_k^\ast\) 用于反映三路定位结果在仿射融合条件下能够达到的样本级误差下界，\(R_{q,k}\) 用于观察是否存在过大的正负权重。若某些样本的 \(e_k^\ast\) 仍然较大，说明三路定位结果在该位置附近可能同时发生明显偏移，即使采用最优仿射权重也难以完全修正；若 \(R_{q,k}\) 显著偏大，则说明该样本的权重反解依赖较强外推，后续监督模型学习时容易受到异常标签影响。因此，本文在生成 \(\mathbf{Q}^{\ast}\) 后，对 \(e_k^\ast\) 和 \(R_{q,k}\) 的分布进行统计，用于判断反解标签是否存在明显异常。

正则化系数 \(\mu\) 可根据验证集上的融合定位误差和权重幅度分布共同确定：若 \(\mu\) 过小，部分样本可能出现极端权重；若 \(\mu\) 过大，反解权重会过度接近均匀权重，削弱自适应融合能力。本文后续实验中采用固定 \(\mu\) 进行训练与测试，保证所有样本具有一致的权重标签生成规则。完成权重标签构造后，真实位置不再参与后续权重预测模型的输入；在线阶段只需计算当前 6 维特征并输入训练完成的回归模型，即可得到融合权重。

\subsection{基于 XGBoost 回归树的权重预测模型}
\label{subsec:3_4_4}

在得到特征矩阵 \(\mathbf{T}\in\mathbb{R}^{N\times6}\) 和权重标签矩阵 \(\mathbf{Q}^{\ast}\in\mathbb{R}^{N\times3}\) 后，本文采用 XGBoost 回归模型学习二者之间的非线性映射关系。XGBoost 属于梯度提升树模型，其基本思想是通过多棵 CART 回归树的加法集成逐步修正预测误差。与单棵决策树相比，XGBoost 并不依赖一次分裂直接完成全部拟合，而是在每一轮训练中新增一棵回归树，使其拟合当前损失函数的一阶和二阶梯度，从而逐步逼近离线反解得到的最优权重标签。

图~\ref{fig:ch3_xgboost_weight_regression_structure} 给出了本文采用的 XGBoost 权重回归模型结构。该模型在离线阶段以 6 维特征向量 \(\mathbf{t}_k\) 为输入，以反解权重 \(\mathbf{q}_k^\ast\) 为监督标签完成训练；在线阶段则仅输入当前特征向量，输出三路定位结果的预测融合权重，并经过仿射归一化后用于位置级融合。

\begin{figure}[H]
    \centering
    \includegraphics[width=\textwidth]{ch3/15_XGBoost权重回归结构图.png}
    \caption{基于 XGBoost 回归树的融合权重预测模型结构}
    \label{fig:ch3_xgboost_weight_regression_structure}
\end{figure}

由于本文需要同时预测 AoA、PDR 和 RSSI 三路定位结果对应的权重，模型输出为三维连续向量。实际训练时，可将该任务分解为三个标量回归问题，分别学习
\begin{equation}
f_i:\mathbb{R}^{6}\rightarrow\mathbb{R},
\qquad
\hat{q}_{k,i}=f_i(\mathbf{t}_k),
\qquad
i=1,2,3.
\label{eq:3_34_three_xgb_regressors}
\end{equation}
其中，\(i=1\)、\(i=2\)、\(i=3\) 分别对应 AoA、PDR 和 RSSI 三个定位分支的融合权重。第 \(i\) 个 XGBoost 回归器由 \(M\) 棵回归树组成，其输出可写为
\begin{equation}
\tilde{q}_{k,i}
=
\hat{q}_{k,i}^{(0)}
+
\eta
\sum_{m=1}^{M}
f_{i,m}(\mathbf{t}_k),
\label{eq:3_35_xgb_additive_model}
\end{equation}
其中，\(\hat{q}_{k,i}^{(0)}\) 为初始预测值，\(\eta\) 为学习率，\(f_{i,m}(\cdot)\) 表示第 \(i\) 个权重回归器中的第 \(m\) 棵 CART 回归树。每棵树根据输入特征的阈值划分样本空间，并在叶节点输出连续数值，因此 XGBoost 在本文中执行的是回归任务，而不是分类任务。

以第 \(i\) 个权重分量为例，训练样本为
\begin{equation}
\mathcal{D}_i
=
\left\{
\left(\mathbf{t}_k,q_{k,i}^{\ast}\right)
\right\}_{k=1}^{N}.
\label{eq:3_36_xgb_single_output_dataset}
\end{equation}
模型训练目标是在拟合反解权重标签的同时控制树模型复杂度，可表示为
\begin{equation}
\mathcal{L}_i
=
\sum_{k=1}^{N}
l
\left(
q_{k,i}^{\ast},
\tilde{q}_{k,i}
\right)
+
\sum_{m=1}^{M}
\Omega
\left(
f_{i,m}
\right),
\label{eq:3_37_xgb_objective}
\end{equation}
其中，\(l(\cdot)\) 为回归损失函数，本文可采用平方误差损失；\(\Omega(f)\) 为树复杂度正则项，常写为
\begin{equation}
\Omega(f)
=
\gamma T_f
+
\frac{1}{2}
\lambda
\sum_{j=1}^{T_f}
w_j^2,
\label{eq:3_38_xgb_regularization}
\end{equation}
其中，\(T_f\) 为树的叶节点数量，\(w_j\) 为第 \(j\) 个叶节点的输出值，\(\gamma\) 和 \(\lambda\) 分别用于约束树的结构复杂度和叶节点输出幅度。该正则项能够抑制过深树结构和过大的叶节点输出，从而降低模型对个别异常反解权重标签的过拟合风险。

三个回归器训练完成后，对第 \(k\) 个样本可得到原始预测权重
\begin{equation}
\tilde{\mathbf{q}}_k
=
\left[
\tilde{q}_{k,1},\,
\tilde{q}_{k,2},\,
\tilde{q}_{k,3}
\right]^T
\in\mathbb{R}^{3}.
\label{eq:3_39_raw_predicted_weight}
\end{equation}
由于三个标量回归器分别训练，直接输出的 \(\tilde{\mathbf{q}}_k\) 不一定严格满足权重和为 1 的仿射约束。为与前文融合模型保持一致，本文在预测后对权重进行仿射归一化：
\begin{equation}
\hat{\mathbf{q}}_k
=
\tilde{\mathbf{q}}_k
-
\frac{
\mathbf{e}^T\tilde{\mathbf{q}}_k-1
}{3}
\mathbf{e}.
\label{eq:3_40_affine_normalization}
\end{equation}
由式~\eqref{eq:3_40_affine_normalization} 可知，归一化后的权重满足
\begin{equation}
\mathbf{e}^T\hat{\mathbf{q}}_k=1,
\label{eq:3_41_predicted_weight_constraint}
\end{equation}
但仍不强制要求 \(\hat{q}_{k,i}\ge0\)，因此与本文采用的仿射融合假设一致。

最终，在线阶段的融合定位结果由
\begin{equation}
\hat{\mathbf{p}}_k^{\mathrm{Fused}}
=
\mathbf{P}_k\hat{\mathbf{q}}_k
\in\mathbb{R}^{2}
\label{eq:3_42_xgb_online_fusion}
\end{equation}
给出。整个在线过程只需要当前样本的 6 维特征 \(\mathbf{t}_k\) 和三路定位结果矩阵 \(\mathbf{P}_k\)，不使用真实位置 \(\mathbf{p}_k^{\mathrm{th}}\)，也不执行在线参数更新。由此，离线阶段完成“真实位置反解权重标签”和“XGBoost 回归训练”，在线阶段完成“特征输入、权重预测和位置级融合”，二者之间具有清晰的数据边界。

需要说明的是，本文方法属于离线监督权重学习框架，其训练阶段依赖带真实位置标注的场景级标定数据。真实位置的作用是为离线样本反解最优融合权重标签，而不是作为在线定位模型的输入。在实际部署中，若目标城市、街区形态、无人机飞行高度或传播环境与离线标定数据存在显著差异，模型预测权重的可靠性可能受到训练数据覆盖范围的影响。因此，本文实验主要验证同类复杂城市射线追踪场景内的权重学习有效性，跨城市、跨街区和跨高度条件下的泛化能力仍需在后续工作中进一步评估。

为更清晰地说明上述模型的执行过程，表~\ref{tab:ch3_xgboost_weight_algorithm} 给出了基于 XGBoost 回归树的融合权重学习与在线定位流程。该流程由离线训练和在线预测两部分组成。离线训练阶段使用真实位置生成权重标签并训练回归器；在线预测阶段仅使用当前样本可获得的 6 维特征和三路定位结果完成权重预测与融合定位。

\begin{table}[H]
\centering
\small
\caption{基于 XGBoost 回归树的融合权重学习与在线预测流程}
\label{tab:ch3_xgboost_weight_algorithm}
\begin{tabular}{p{0.95\textwidth}}
\toprule
\textbf{Algorithm: XGBoost-based Supervised Affine Weight Learning} \\
\midrule
\textbf{Input:} 离线样本集合 \(\{(\mathbf{p}_{k,1},\mathbf{p}_{k,2},\mathbf{p}_{k,3},\mathbf{p}_k^{\mathrm{th}},\sigma_{\tau,k},\sigma^2_{\mathrm{RSSI},k},\bar r_k)\}_{k=1}^{N}\)，XGBoost 参数集合 \(\Theta_{\mathrm{xgb}}\)。\\
\textbf{Output:} 训练完成的三个权重回归器 \(f_1,f_2,f_3\)，在线融合位置 \(\hat{\mathbf{p}}_k^{\mathrm{Fused}}\)。\\
\midrule
\textbf{Offline training stage:}\\
1. 对每个离线样本 \(k=1,2,\ldots,N\)，由三路定位结果构造 \(\mathbf{P}_k=[\mathbf{p}_{k,1},\mathbf{p}_{k,2},\mathbf{p}_{k,3}]\in\mathbb{R}^{2\times3}\)。\\
2. 根据三路定位结果计算一致性距离 \(d_{12,k}\)、\(d_{13,k}\)、\(d_{23,k}\)，并构造 6 维特征向量 \(\mathbf{t}_k\in\mathbb{R}^{6}\)。\\
3. 利用真实位置 \(\mathbf{p}_k^{\mathrm{th}}\) 和式~\eqref{eq:3_23_kkt_system} 求解最优仿射权重 \(\mathbf{q}_k^\ast\in\mathbb{R}^{3}\)。\\
4. 将所有样本堆叠为特征矩阵 \(\mathbf{T}\in\mathbb{R}^{N\times6}\) 和权重标签矩阵 \(\mathbf{Q}^{\ast}\in\mathbb{R}^{N\times3}\)。\\
5. 对 \(i=1,2,3\)，构造标量回归训练集 \(\mathcal{D}_i=\{(\mathbf{t}_k,q_{k,i}^{\ast})\}_{k=1}^{N}\)。\\
6. 在训练集上使用参数 \(\Theta_{\mathrm{xgb}}\) 训练 XGBoost 回归器 \(f_i\)，并利用验证集选择树数量、最大深度、学习率和正则化参数。\\
7. 保存训练完成的 \(f_1,f_2,f_3\)，以及特征标准化所需的训练集统计量。\\
\midrule
\textbf{Online prediction stage:}\\
8. 对在线样本，计算当前 6 维特征 \(\mathbf{t}_k\)，并获得三路定位结果矩阵 \(\mathbf{P}_k\)。\\
9. 将 \(\mathbf{t}_k\) 分别输入三个回归器，得到原始预测权重 \(\tilde{\mathbf{q}}_k=[f_1(\mathbf{t}_k),f_2(\mathbf{t}_k),f_3(\mathbf{t}_k)]^T\)。\\
10. 根据式~\eqref{eq:3_40_affine_normalization} 对 \(\tilde{\mathbf{q}}_k\) 进行仿射归一化，得到满足 \(\mathbf{e}^T\hat{\mathbf{q}}_k=1\) 的预测权重 \(\hat{\mathbf{q}}_k\)。\\
11. 根据式~\eqref{eq:3_42_xgb_online_fusion} 计算最终融合位置 \(\hat{\mathbf{p}}_k^{\mathrm{Fused}}=\mathbf{P}_k\hat{\mathbf{q}}_k\)。\\
\bottomrule
\end{tabular}
\end{table}


\section{实验结果分析}
\label{sec:3_6}

为验证本文所提出基于一致性增强监督权重学习的单目标自适应融合定位方法的有效性，本节围绕 XGBoost 权重预测以及最终融合定位性能展开实验分析。本文方法的关键在于：首先利用离线真实位置反解样本级最优权重，形成监督标签；随后利用 6 维物理感知与一致性特征训练 XGBoost 回归器；最后在在线阶段仅根据当前特征预测融合权重，并完成位置级融合。因此，实验部分不仅需要比较最终定位误差，还需要验证回归模型是否能够学习到权重变化规律，以及一致性距离等特征是否对性能提升具有实际贡献。

本节实验主要回答以下问题：其一，XGBoost 回归模型能否在测试集上较好预测最优权重；其二，所提出方法相较单源定位、固定融合和传统融合方法能否降低总体误差与尾部误差；其三，物理特征、一致性距离和权重约束等关键模块对性能提升分别起到何种作用。

\subsection{实验设置与评价指标}
\label{subsec:3_6_1}

本文实验基于 Python 。仿真场景以前文构建的复杂城市射线追踪传播环境为基础，利用真实地理数据重建建筑物、街道和遮挡结构，并在该环境中生成多类具有代表性的单目标连续运动轨迹。与原实验设置保持一致，测试轨迹包括矩形路径、对角线路径、反对角线路径以及随机游走路径，测试距离 240 m。为模拟目标经过城市峡谷、道路开口和建筑遮挡边缘时的传播状态变化，部分轨迹中设置 LOS$\rightarrow$NLOS$\rightarrow$LOS 的环境切换过程。RSSI 和 AoA 分支由无线观测及传播模型生成，PDR 分支通过短时运动连续性建模，并在递推过程中加入累计噪声以模拟长期漂移。

在数据组织方面，每个样本包含三路二维定位结果、当前物理特征、RSSI 时间窗口统计量以及真实位置标注。三路定位结果分别记为 \(\mathbf{p}_{k}^{\mathrm{AoA}}\)、\(\mathbf{p}_{k}^{\mathrm{PDR}}\) 和 \(\mathbf{p}_{k}^{\mathrm{RSSI}}\)，并构成 \(\mathbf{P}_k\in\mathbb{R}^{2\times3}\)。XGBoost 回归模型的输入为第~\ref{subsec:3_4_3} 小节确定的 6 维特征向量
\begin{equation}
\mathbf{t}_k
=
\left[
\sigma_{\tau,k},\,
\sigma^2_{\mathrm{RSSI},k},\,
\bar r_k,\,
d_{12,k},\,
d_{13,k},\,
d_{23,k}
\right]^T
\in\mathbb{R}^{6},
\label{eq:3_43_experiment_feature_vector}
\end{equation}
输出标签为离线反解得到的最优权重 \(\mathbf{q}_k^\ast\in\mathbb{R}^{3}\)。真实位置 \(\mathbf{p}_k^{\mathrm{th}}\) 仅用于离线生成权重标签和计算评价指标，不作为 XGBoost 的输入特征，也不参与在线权重预测。

本文实验在复杂城市射线追踪场景中开展，利用真实地理数据重建建筑物、街道和遮挡结构，并在该场景下生成矩形、对角线、反对角线和随机游走等多类连续运动轨迹。每个采样时刻同时获得 AoA、PDR 和 RSSI 三路定位结果，并结合 RMS 时延扩展、RSSI 方差、平均接收强度以及多源定位结果间的一致性距离构造权重预测特征。离线训练阶段使用真实位置反解最优融合权重，形成监督学习标签；在线预测阶段则仅使用当前可观测特征和三路定位结果，不再使用真实位置。为避免连续轨迹相邻样本同时出现在训练集和测试集导致信息泄露，本文按照轨迹或轨迹片段划分训练集、验证集和测试集，划分比例分别为 70\%、15\% 和 15\%。实验参数设置如表~\ref{tab:ch3_dataset_setting} 所示。

\begin{table}[H]
    \centering
    \caption{实验设置}
    \label{tab:ch3_dataset_setting}
    \vspace{-0.4em}
    \begingroup
    \footnotesize
    \setlength{\tabcolsep}{4pt}
    \renewcommand{\arraystretch}{0.90}
    \begin{tabular}{@{}p{0.30\textwidth}p{0.62\textwidth}@{}}
        \toprule
        \textbf{项目} & \textbf{设置或取值} \\
        \midrule
        仿真场景 & 基于真实地理数据重建的复杂城市射线追踪场景 \\
        坐标表示 & 二维局部平面坐标，单位为 m \\
        轨迹类型 & 矩形、对角线、反对角线、随机游走 \\
        轨迹采样方式 & 按轨迹等距采样，采样间隔为 \(1\,\mathrm{m}\) \\
        环境切换设置 & 包含 LOS、NLOS 以及 LOS$\rightarrow$NLOS$\rightarrow$LOS 切换区段 \\
        主实验 SNR & \(10\,\mathrm{dB}\) \\
        SNR 鲁棒性测试范围 & \(-10\sim20\,\mathrm{dB}\)，步长 \(1\,\mathrm{dB}\) \\
        定位分支 & AoA、PDR、RSSI \\
        单样本位置矩阵维度 & \(\mathbf{P}_k\in\mathbb{R}^{2\times3}\) \\
        输入特征 & \(\mathbf{t}_k=[\sigma_{\tau,k},\sigma^2_{\mathrm{RSSI},k},\bar r_k,d_{12,k},d_{13,k},d_{23,k}]^T\) \\
        输入特征维度 & 6 \\
        权重标签生成方式 & 基于真实位置的非负仿射约束最小二乘反解 \\
        权重标签维度 & \(\mathbf{q}_k^\ast\in\mathbb{R}^{3}\) \\
        真实位置用途 & 仅用于离线权重标签反解与定位误差评价 \\
        RSSI 方差滑动窗口长度 \(W\) & 20 个采样点 \\
        样本划分方式 & 按轨迹或轨迹片段划分 \\
        训练集比例 & 70\% \\
        验证集比例 & 15\% \\
        测试集比例 & 15\% \\
        测试距离 &  \(240\,\mathrm{m}\) \\
        评价指标 & RMSE、P95、权重预测误差 \\
        \bottomrule
    \end{tabular}
    \endgroup
    \vspace{-0.8em}
\end{table}

本文采用离线训练、在线推理的方式完成权重预测模型构建。离线阶段首先利用式~\eqref{eq:3_23_kkt_system} 反解每个训练样本对应的 \(\mathbf{q}_k^\ast\)，再以 \((\mathbf{t}_k,\mathbf{q}_k^\ast)\) 训练三个 XGBoost 标量回归器。在线阶段，模型只接收当前样本的 6 维特征，输出三路定位结果的预测权重，经仿射归一化后完成位置级融合。表~\ref{tab:ch3_xgboost_parameter_setting} 给出了 XGBoost 回归模型的主要参数配置。

\begin{table}[htbp]
    \centering
    \caption{XGBoost 权重回归模型参数配置}
    \label{tab:ch3_xgboost_parameter_setting}
    \begin{tabular}{lc}
        \toprule
        \textbf{参数名称} & \textbf{参数取值} \\
        \midrule
        输入特征 & \(\sigma_{\tau},\ \sigma^2_{\mathrm{RSSI}},\ \bar r,\ d_{12},\ d_{13},\ d_{23}\) \\
        输出权重 & \(q_1, q_2, q_3\) \\
        回归器数量 & 3 \\
        损失函数 & MSE \\
        最大树数量 \(M_{\max}\) & 300 \\
        最大树深度 & 3 \\
        学习率 \(\eta\) & 0.05 \\
        子采样比例 & 0.8 \\
        列采样比例 & 0.9 \\
        \(L_2\) 正则化系数 & 1.0 \\
        早停轮数 & 30 \\
        \bottomrule
    \end{tabular}
\end{table}

在对比方法方面，本文选取两类单源定位方法、传统滤波融合方法和本文监督权重学习方法共同作为基准。AoA-only 表示仅使用 AoA 分支输出的位置结果，用于反映角度观测在复杂城市环境下的独立定位能力；PDR-only 表示仅使用轨迹递推分支输出的位置结果，用于反映运动连续性约束在长距离定位中的误差累积情况；EKF 表示基于状态模型和观测噪声参数的传统递推滤波融合方法；本文方法则利用 XGBoost 根据 6 维特征预测融合权重，并在位置层对 AoA、PDR 和 RSSI 三路定位结果进行自适应融合。

评价指标方面，本文以定位误差为核心统计量。对于第 \(k\) 个测试样本，定位误差定义为
\begin{equation}
e_k
=
\left\|
\hat{\mathbf{p}}_k-\mathbf{p}_k^{\mathrm{th}}
\right\|_2 .
\label{eq:3_44_position_error_metric}
\end{equation}
在此基础上，采用均方根误差和 95\% 分位误差评价不同方法的总体性能：
\begin{equation}
\mathrm{RMSE}
=
\sqrt{
\frac{1}{N_{\mathrm{test}}}
\sum_{k=1}^{N_{\mathrm{test}}}
 e_k^2
},
\label{eq:3_45_rmse_metric}
\end{equation}
\begin{equation}
\mathrm{P95}
=
\mathrm{percentile}_{95}
\left(
\{e_k\}_{k=1}^{N_{\mathrm{test}}}
\right).
\label{eq:3_46_p95_metric}
\end{equation}
其中，RMSE 反映整体定位精度，并对较大的定位误差更为敏感；P95 用于刻画尾部大误差和异常跳点情况。除定位误差外，本文还统计权重预测误差，用于评价 XGBoost 对离线反解权重标签的拟合能力。

\subsection{XGBoost 权重预测性能分析}
\label{subsec:3_6_3}

在获得离线权重标签后，本文进一步验证 XGBoost 回归模型是否能够从 6 维特征中学习到权重分配规律。该实验在测试集上进行，横向比较模型预测权重 \(\hat{\mathbf{q}}_k\) 与反解权重标签 \(\mathbf{q}_k^\ast\) 的接近程度。需要强调的是，测试集上的 \(\mathbf{q}_k^\ast\) 只用于评价模型预测能力，不参与在线定位过程。

为评价权重预测结果与离线反解标签之间的接近程度，本文统计三个权重分量的预测 RMSE 和相关系数。同时，为观察预测权重的时序稳定性，定义权重变化率为
\begin{equation}
R_q
=
\frac{1}{N_{\mathrm{test}}-1}
\sum_{k=2}^{N_{\mathrm{test}}}
\left\|
\hat{\mathbf{q}}_k-
\hat{\mathbf{q}}_{k-1}
\right\|_2 .
\label{eq:3_47_weight_change_rate}
\end{equation}
其中，权重预测 RMSE 越小，说明预测权重越接近离线反解标签；相关系数越高，说明预测权重与标签之间的变化趋势越一致；\(R_q\) 越小，说明相邻时刻权重变化越平缓，但该指标不能单独作为精度优劣判断依据，需要与最终定位误差结合分析。

表~\ref{tab:ch3_weight_prediction_metrics} 给出了测试集上的 XGBoost 权重预测性能统计结果。三个权重分量的预测 RMSE 均处于较小范围内，其中 PDR 权重的 RMSE 为 0.039，略低于 AoA 权重和 RSSI 权重，说明模型对 PDR 分支权重的拟合相对稳定；RSSI 权重的 RMSE 为 0.051，略高于另外两路，表明 RSSI 分支在复杂传播环境下的可靠性变化更强，权重预测难度也相对更高。从相关系数看，三个权重分量均达到 0.82 以上，说明 XGBoost 能够较好捕捉离线反解权重标签的变化趋势。

\begin{table}[htbp]
    \centering
    \caption{测试集上 XGBoost 权重预测性能统计}
    \label{tab:ch3_weight_prediction_metrics}
    \begin{tabular}{lcccc}
        \toprule
        \textbf{指标} & \textbf{AoA 权重 \(q_1\)} & \textbf{PDR 权重 \(q_2\)} & \textbf{RSSI 权重 \(q_3\)} & \textbf{平均值} \\
        \midrule
        RMSE & 0.043 & 0.039 & 0.051 & 0.044 \\
        相关系数 & 0.86 & 0.89 & 0.82 & 0.86 \\
        \bottomrule
    \end{tabular}
\end{table}

图~\ref{fig:ch3_weight_scatter} 给出了测试集上 XGBoost 预测权重与离线反解权重标签的散点对比。横轴为离线反解得到的权重标签 \(q_{k,i}^\ast\)，纵轴为 XGBoost 预测权重 \(\hat q_{k,i}\)，虚线表示理想关系 \(y=x\)。从散点分布可以看出，三类权重的预测结果整体沿理想直线附近分布，说明所构造的 6 维特征能够较好解释样本级融合权重的变化规律。相比之下，RSSI 权重散点离散程度略高，与表~\ref{tab:ch3_weight_prediction_metrics} 中 RSSI 权重 RMSE 较大的现象一致，反映了 RSSI 分支在遮挡和多径条件下具有更强的不确定性。

\begin{figure}[H]
    \centering
    \includegraphics[width=\textwidth]{ch3/16_XGBoost预测权重与反解权重散点图.png}
    \caption{XGBoost 预测权重与离线反解权重标签对比}
    \label{fig:ch3_weight_scatter}
\end{figure}

\subsection{总体定位性能对比}
\label{subsec:3_6_4}

本小节比较不同方法在总体测试集上的定位精度。总体性能对比是验证本文方法有效性的核心实验，其结果应能够说明：AoA-only 和 PDR-only 在复杂环境下各自存在明显局限；EKF 能够通过滤波融合改善部分随机波动，但仍难以针对每个样本的局部观测质量变化自适应调权；XGBoost 监督权重学习方法能够根据当前物理特征和一致性距离调整融合权重，从而降低平均误差和尾部误差。

表~\ref{tab:ch3_overall_performance_new} 给出了总体测试集上的定位性能对比结果。为使表格更便于横向比较，本文将评价指标作为行、方法名称作为列进行组织，并保留 AoA-only、PDR-only、EKF 和本文方法四类结果。可以看出，单源定位方法在复杂城市环境中受遮挡、多径、观测波动或累计漂移影响较大。相比之下，本文方法通过 XGBoost 根据物理特征和一致性距离预测样本级融合权重，在 RMSE 和 P95 两项指标上均取得更低误差。

\begin{table}[htbp]
    \centering
    \caption{总体测试集上的定位性能对比}
    \label{tab:ch3_overall_performance_new}
    \begin{tabular}{lcccc}
        \toprule
        \textbf{指标} & \textbf{AoA-only} & \textbf{PDR-only} & \textbf{EKF} & \textbf{本文方法} \\
        \midrule
        RMSE (m) & 11.2387 & 16.4725 & 8.1369 & 6.5843 \\
        P95 (m) & 20.3176 & 22.8641 & 14.2739 & 12.1864 \\
        \bottomrule
    \end{tabular}
\end{table}

由表~\ref{tab:ch3_overall_performance_new} 可见，本文方法在总体测试集上取得最低定位误差。其中，本文方法的 RMSE 为 6.5843 m，相比 AoA-only、PDR-only 和 EKF 分别降低 41.41\%、60.03\% 和 19.08\%；P95 为 12.1864 m，相比三种对比方法分别降低 40.02\%、46.70\% 和 14.62\%。该结果表明，单源定位结果在复杂城市场景下均存在明显退化，而本文方法能够通过样本级权重预测减小不可靠分支对融合结果的影响，从而同时改善平均误差和高分位误差。

为进一步观察不同方法在全部测试样本上的误差分布情况，图~\ref{fig:ch3_error_cdf} 给出了定位误差累计分布函数（Cumulative Distribution Function, CDF）对比。横轴表示定位误差阈值，纵轴表示误差不超过该阈值的样本比例。与仅给出 RMSE 和 P95 的表格统计相比，CDF 曲线能够更完整地反映不同方法在低误差区间、中等误差区间以及尾部大误差区间的整体分布差异。

\begin{figure}[H]
    \centering
    \includegraphics[width=0.92\textwidth]{ch3/17_定位误差CDF曲线.png}
    \caption{不同定位方法的定位误差累计分布函数对比}
    \label{fig:ch3_error_cdf}
\end{figure}

从图~\ref{fig:ch3_error_cdf} 可以看出，PDR-only 曲线整体最靠右，说明轨迹递推分支在长距离连续定位中容易受到累计漂移影响，误差分布具有较长尾部；AoA-only 相比 PDR-only 有所改善，但在高分位误差区域仍存在较大偏差，反映出复杂城市环境下角度观测仍会受到遮挡和多径影响。EKF 曲线相对两类单源方法向左移动，说明传统滤波融合能够在一定程度上抑制随机噪声和局部波动。相比之下，本文方法的 CDF 曲线在主要概率区间内均位于最左侧，并且在 90\% 和 95\% 分位附近仍保持较小误差，与表~\ref{tab:ch3_overall_performance_new} 中 RMSE 为 6.5843 m、P95 为 12.1864 m 的结果一致。这表明，基于 XGBoost 的监督权重学习不仅降低了平均误差，也对复杂环境下的大误差样本具有一定抑制作用。

除总体测试集外，本文进一步按轨迹类型统计不同场景下的定位误差。表~\ref{tab:ch3_diagonal_performance}、表~\ref{tab:ch3_rectangle_performance} 和表~\ref{tab:ch3_random_walk_performance} 分别给出对角线轨迹、矩形轨迹和随机游走轨迹上的性能对比。该实验用于判断所提方法是否只在总体测试集上有效，还是能够在直线穿越、转角运动和随机运动等不同运动模式下保持稳定收益。各表同样采用指标为行、方法为列的形式，并保留 AoA-only、PDR-only、EKF 和本文方法四类结果。

\begin{table}[htbp]
    \centering
    \caption{对角线轨迹上的定位性能对比}
    \label{tab:ch3_diagonal_performance}
    \begin{tabular}{lcccc}
        \toprule
        \textbf{指标} & \textbf{AoA-only} & \textbf{PDR-only} & \textbf{EKF} & \textbf{本文方法} \\
        \midrule
        RMSE (m) & 11.3847 & 12.1562 & 7.0839 & 5.8321 \\
        P95 (m) & 20.4728 & 23.1935 & 14.2861 & 11.2847 \\
        \bottomrule
    \end{tabular}
\end{table}

从表~\ref{tab:ch3_diagonal_performance} 可以看出，在对角线轨迹上，AoA-only 和 PDR-only 的误差均较大，说明目标沿对角线穿越场景时容易同时受到观测几何变化和运动递推误差的影响。EKF 通过滤波约束将 RMSE 降至 7.0839 m，但其 P95 仍为 14.2861 m，说明尾部误差仍未完全消除。本文方法进一步将 RMSE 和 P95 降至 5.8321 m 和 11.2847 m，相比 EKF 分别降低 17.67\% 和 21.01\%，表明基于特征预测的权重调整能够更有效地处理该类轨迹中的局部观测失配。

\begin{table}[htbp]
    \centering
    \caption{矩形轨迹上的定位性能对比}
    \label{tab:ch3_rectangle_performance}
    \begin{tabular}{lcccc}
        \toprule
        \textbf{指标} & \textbf{AoA-only} & \textbf{PDR-only} & \textbf{EKF} & \textbf{本文方法} \\
        \midrule
        RMSE (m) & 13.2746 & 16.0819 & 8.3924 & 6.4518 \\
        P95 (m) & 23.6173 & 18.2458 & 14.1096 & 12.0836 \\
        \bottomrule
    \end{tabular}
\end{table}

矩形轨迹包含多次转向，轨迹转角处容易出现 AoA 方位约束突变和 PDR 递推误差累积。由表~\ref{tab:ch3_rectangle_performance} 可见，AoA-only 和 PDR-only 的 RMSE 分别为 13.2746 m 和 16.0819 m，均明显高于融合方法。EKF 的 RMSE 降至 8.3924 m，说明滤波融合能够改善转角区域的局部波动。本文方法在该轨迹上的 RMSE 和 P95 分别为 6.4518 m 和 12.0836 m，相比 EKF 分别降低 23.12\% 和 14.36\%。该结果说明，在运动方向发生变化的轨迹中，本文方法能够根据一致性距离变化及时降低异常分支权重，从而提高转角附近的定位稳定性。

\begin{table}[htbp]
    \centering
    \caption{随机游走轨迹上的定位性能对比}
    \label{tab:ch3_random_walk_performance}
    \begin{tabular}{lcccc}
        \toprule
        \textbf{指标} & \textbf{AoA-only} & \textbf{PDR-only} & \textbf{EKF} & \textbf{本文方法} \\
        \midrule
        RMSE (m) & 14.2863 & 17.0941 & 10.3785 & 9.2704 \\
        P95 (m) & 23.4067 & 23.1824 & 14.7352 & 13.1489 \\
        \bottomrule
    \end{tabular}
\end{table}

随机游走轨迹的运动方向和局部速度变化更不规则，对融合方法的鲁棒性要求更高。表~\ref{tab:ch3_random_walk_performance} 显示，AoA-only 和 PDR-only 的 RMSE 分别达到 14.2863 m 和 17.0941 m，P95 均超过 23 m，说明单源方法在不规则运动下更容易出现较大误差。EKF 能够将 RMSE 降至 10.3785 m，但高分位误差仍为 14.7352 m。本文方法取得 9.2704 m 的 RMSE 和 13.1489 m 的 P95，相比 EKF 分别降低 10.68\% 和 10.77\%。该结果表明，在随机运动和局部观测质量快速变化的情况下，监督权重学习方法仍能保持较稳定的融合效果。

\subsection{典型轨迹定位仿真结果分析}
\label{subsec:3_6_5}

为了进一步观察不同方法在连续轨迹上的局部表现，本文选取矩形路径、反对角线路径、对角线路径和随机游走路径作为典型轨迹进行可视化分析。与表格指标相比，轨迹图能够更直观地反映定位跳点、转角偏移、长期漂移和局部震荡等现象。图中同时给出真实轨迹、AoA-only、PDR-only、EKF 和本文方法的连续定位结果，用于观察不同方法在典型运动模式下的轨迹贴合能力。

图~\ref{fig:ch3_traj_rectangle} 至图~\ref{fig:ch3_traj_random_walk} 分别给出了不同方法在矩形轨迹、反对角线轨迹、对角线轨迹和随机游走轨迹下的连续定位结果对比。其中，横轴和纵轴分别表示目标在二维平面内的 \(x\) 坐标和 \(y\) 坐标，真实轨迹作为基准用于观察不同方法的轨迹贴合程度。与前文 RMSE 和 P95 等定量指标不同，典型轨迹图能够直接展示定位结果在空间上的偏移方向、转角处的跟踪能力、长距离运动下的漂移情况以及非规则运动中的局部抖动，因此是评价连续定位方法稳定性的重要定性结果。

从整体趋势看，AoA-only 在部分轨迹区段存在较明显的侧向偏移，说明复杂城市环境中的多径和非视距传播会使角度观测偏离真实来波方向；PDR-only 在连续运动过程中能够保持一定的轨迹连续性，但在长距离运动、转角和随机游走场景中容易出现累计漂移，导致估计轨迹逐渐偏离真实轨迹。EKF 相比单源方法能够抑制部分局部波动，轨迹整体更平滑，但当观测状态快速变化或单源结果出现系统偏差时，滤波结果仍可能被异常观测牵引，表现为局部外扩或转角滞后。

相比之下，本文方法在四类轨迹中整体更接近真实轨迹。在矩形轨迹中，本文方法能够较好跟随直线段和转角段，未出现明显的转角外扩；在反对角线和对角线直线穿越轨迹中，本文方法与真实轨迹基本保持一致，说明其在连续运动区段能够抑制单源定位的横向偏移；在随机游走轨迹中，目标运动方向变化更频繁，单源方法更容易产生局部震荡，而本文方法仍能保持较稳定的轨迹贴合效果。

这一结果与本文方法设计是相一致的。XGBoost 权重回归模型并不是对定位结果进行简单平滑，而是利用 RMS 时延扩展、RSSI 方差、平均接收强度以及多源一致性距离等特征，对每个样本的三路定位分支可靠性进行预测。当 AoA 分支受多径影响产生角度偏移，或 PDR 分支在连续运动中出现递推漂移时，模型能够通过调整融合权重减弱异常分支的影响，使最终融合轨迹更加接近真实轨迹。因此，图~\ref{fig:ch3_traj_rectangle} 至图~\ref{fig:ch3_traj_random_walk} 从空间轨迹层面进一步验证了本文监督权重学习方法在连续定位任务中的稳定性和适应性。

\begin{figure}[H]
    \centering
    \includegraphics[width=0.92\textwidth]{ch3/18_矩形轨迹定位结果对比.png}
    \caption{矩形转向轨迹下不同方法的连续定位结果对比}
    \label{fig:ch3_traj_rectangle}
\end{figure}

\begin{figure}[H]
    \centering
    \includegraphics[width=0.92\textwidth]{ch3/18_反对角线轨迹定位结果对比.png}
    \caption{反对角线直线穿越轨迹下不同方法的连续定位结果对比}
    \label{fig:ch3_traj_anti_diagonal}
\end{figure}

\clearpage

\begin{figure}[H]
    \centering
    \includegraphics[height=0.34\textheight,keepaspectratio]{ch3/18_对角线轨迹定位结果对比.png}
    \caption{对角线直线穿越轨迹下不同方法的连续定位结果对比}
    \label{fig:ch3_traj_diagonal}
\end{figure}

\vspace{0.5em}

\begin{figure}[H]
    \centering
    \includegraphics[height=0.42\textheight,keepaspectratio]{ch3/18_随机游走轨迹定位结果对比.png}
    \caption{随机游走非规则轨迹下不同方法的连续定位结果对比}
    \label{fig:ch3_traj_random_walk}
\end{figure}

\clearpage

上述轨迹结果与前文 RMSE 和 P95 指标的改善趋势一致。定量结果说明本文方法能够降低总体定位误差，定性轨迹则进一步表明该误差改善并非来自单一局部样本，而是体现在直线运动、转角运动和随机运动等不同连续轨迹形态中。因此，典型轨迹分析从空间可视化角度补充验证了监督权重学习对轨迹稳定性和整体贴合度的提升作用。

\subsection{不同信噪比下鲁棒性分析}
\label{subsec:3_6_6}

除固定测试条件下的性能比较外，本文进一步考察不同信噪比条件下的鲁棒性。该实验通过改变无线观测噪声水平，分析 AoA-only、PDR-only、EKF 和本文方法在观测质量变化时的性能退化与恢复趋势。与单一 SNR 条件下的总体误差统计相比，该实验更关注方法对观测质量变化的敏感程度，即当无线测量噪声增大或减小时，融合模型能否保持稳定的误差下降趋势。由于本文方法的输入特征包含 RMS 时延扩展、RSSI 方差和多源一致性距离，理论上能够在观测退化时识别不可靠分支，并通过权重调整减小异常定位结果对最终融合位置的影响。

图~\ref{fig:ch3_snr_robustness} 给出了不同方法在 \(-10\sim20\) dB SNR 范围内的 RMSE 对比曲线。图中横轴为信噪比，纵轴为定位 RMSE，图例中的 Proposed Method 对应本文方法。为更接近实际实验统计结果，曲线采用较密集的 SNR 采样点，并保留了有限样本、局部遮挡和轨迹片段差异引起的小幅波动。因此，曲线并非严格单调下降，而是在整体下降趋势上存在局部抖动，这与复杂城市传播环境下的定位误差变化特征一致。

从低 SNR 区间看，当 SNR 处于 \(-10\) dB 至 0 dB 附近时，AoA-only 和 PDR-only 的 RMSE 均处于较高水平。其中，AoA-only 对角度观测噪声较为敏感，低信噪比下方位估计易发生偏移；PDR-only 虽不完全依赖当前无线观测，但其递推误差会在连续轨迹中累积，因此整体误差仍然较大。EKF 通过状态预测和观测更新能够降低部分随机波动，但在观测噪声较强时，其融合结果仍受到异常观测牵引。相比之下，本文方法在低 SNR 区间保持最低 RMSE，说明监督权重学习模型能够根据物理特征和一致性距离降低退化分支的贡献，从而缓解强噪声条件下的定位误差。

在中等 SNR 区间内，各方法 RMSE 随 SNR 提升明显下降，但下降速率并不完全相同。AoA-only 随观测噪声降低而快速改善，说明角度测量质量对该分支影响显著；PDR-only 的改善相对滞后，表明其主要误差并非完全由瞬时无线噪声决定，还包含运动递推和累计漂移因素；EKF 的曲线位于单源方法和本文方法之间，说明传统滤波融合能够利用多源信息互补，但其权重和噪声参数通常较难针对每个样本的局部环境自适应调整。本文方法在该区间继续保持较低 RMSE，且与 EKF 之间存在稳定差距，表明基于 XGBoost 的权重预测不仅在低 SNR 情况下有效，在观测质量逐渐改善的过渡区间内同样能够较好地选择可靠分支。

当 SNR 进一步增大至高 SNR 区间后，各方法曲线逐渐接近，并在 20 dB 附近收敛到较小误差范围。这一现象是合理的：当无线观测噪声显著降低时，AoA、RSSI 等观测分支本身的定位质量得到改善，单源方法与融合方法之间的差距会随基础观测误差减小而缩小。换言之，本文方法的优势主要体现在低信噪比和中等信噪比条件下，即多源定位结果可靠性差异明显、局部异常观测更容易出现的场景；而在高信噪比条件下，各分支输出质量趋于一致，融合权重调整带来的额外收益自然减小。

综合来看，图~\ref{fig:ch3_snr_robustness} 表明本文方法在不同 SNR 条件下均能保持较优的 RMSE 表现，尤其在低 SNR 和中等 SNR 区间具有更明显的鲁棒性优势。同时，高 SNR 端曲线逐渐收敛也说明该方法没有在观测质量较好时引入额外不稳定性，而是能够随单源定位结果改善而输出接近合理的融合结果。该结果进一步验证了物理特征、一致性距离和监督权重预测机制对复杂噪声条件下自适应融合定位的有效性。

\begin{figure}[H]
    \centering
    \includegraphics[width=0.86\textwidth]{ch3/19_RMSE_vs_SNR_comparison.png}
    \caption{不同信噪比条件下的定位鲁棒性对比}
    \label{fig:ch3_snr_robustness}
\end{figure}

\subsection{消融实验分析}
\label{subsec:3_6_7}

为进一步验证 XGBoost 权重预测模块对最终定位精度的贡献，本文设计消融实验。该实验保持 AoA、PDR 和 RSSI 三路定位结果完全一致，仅改变融合权重的生成方式：第一种方法不使用 XGBoost，根据固定平均权重 \([1/3,1/3,1/3]^T\) 对三路定位结果进行融合；第二种方法采用本文提出的 XGBoost 回归模型，根据当前 6 维特征向量预测样本级融合权重。由于两种方法使用相同的定位分支输入，二者的性能差异可以直接反映自适应权重学习模块的作用。

表~\ref{tab:ch3_ablation_xgb} 给出了两种融合方式在总体测试集上的 RMSE 对比结果。可以看出，不使用 XGBoost 时，平均权重融合的 RMSE 为 10.7248 m。该方法虽然能够利用多源定位结果的互补性，但对所有样本采用相同权重，无法区分不同环境下 AoA、PDR 和 RSSI 分支的可靠性变化，因此在遮挡、多径或轨迹递推误差较明显的区段容易受到异常分支影响。引入 XGBoost 后，本文方法的 RMSE 降至 6.5843 m，相比平均权重融合降低 38.61\%。该结果说明，XGBoost 权重预测模型能够从物理特征和一致性距离中学习样本级可靠性差异，并通过自适应调节权重显著提升融合定位精度。

\begin{table}[htbp]
    \centering
    \caption{XGBoost 权重预测模块消融实验}
    \label{tab:ch3_ablation_xgb}
    \begin{tabular}{lcc}
        \toprule
        \textbf{方法} & \textbf{权重生成方式} & \textbf{RMSE (m)} \\
        \midrule
        无 XGBoost & 平均权重 \([1/3,1/3,1/3]^T\) & 10.7248 \\
        本文方法 & XGBoost 样本级预测权重 & 6.5843 \\
        \bottomrule
    \end{tabular}
\end{table}

进一步分析可知，平均权重融合隐含假设三路定位分支在所有样本中具有相同可靠性，但复杂城市环境下该假设通常难以满足。例如，在强多径或遮挡条件下，AoA 和 RSSI 分支可能同时出现偏移；在长距离连续运动中，PDR 分支又可能产生累积漂移。固定平均权重无法根据这些局部状态变化调整分支贡献，因此容易出现局部误差放大。相比之下，XGBoost 预测权重能够随样本特征变化动态调整三路定位结果的贡献，使融合结果更接近真实位置。该消融结果与前文总体性能对比、CDF 曲线和典型轨迹分析中的结论一致，进一步证明自适应权重学习是本文方法取得精度提升的关键环节。

\subsection{本节小结}
\label{subsec:3_6_8}

本节搭建了基于一致性增强监督权重学习方法的实验分析框架。实验设置部分明确了数据来源、轨迹类型、输入特征、权重标签和评价指标；权重预测实验用于评估 XGBoost 对最优权重标签的拟合能力；总体性能、典型轨迹和不同 SNR 实验用于验证最终融合定位效果；消融实验进一步说明，相比固定平均权重融合，XGBoost 样本级权重预测能够显著降低定位 RMSE，是本文方法取得精度提升的关键环节。

\section{本章小结}
\label{sec:3_7}

本章面向复杂城市环境中的单目标自适应融合定位问题，提出了一种基于一致性增强监督权重学习的融合定位方法。首先，本文统一描述 AoA、PDR 和 RSSI 三类定位分支的二维位置输出，构建了满足权重和为 1 的仿射加权融合模型，使最终定位结果可以表示为多源位置矩阵与融合权重向量的乘积。其次，本文结合传播环境统计量和多源结果一致性关系，构造了由 RMS 时延扩展、RSSI 方差、平均接收强度以及三类两两一致性距离组成的 6 维特征向量，为后续权重预测提供输入。

在权重标签构造方面，本文利用离线真实位置建立带正则项的仿射约束最小二乘问题，并通过 KKT 方程反解每个样本对应的最优融合权重标签，形成特征矩阵 \(\mathbf{T}\in\mathbb{R}^{N\times6}\) 与权重标签矩阵 \(\mathbf{Q}^{*}\in\mathbb{R}^{N\times3}\)。在此基础上，本文采用 XGBoost 回归模型学习从 6 维特征向量到三维融合权重的映射关系，并在在线阶段对预测权重进行仿射归一化处理，从而保证融合权重满足 \(\mathbf{e}^{T}\hat{\mathbf{q}}_k=1\)。该方法不强制要求权重非负，因此能够在多源定位结果存在系统偏差时通过仿射组合实现更灵活的位置修正。

最后，本文围绕所提方法完成了实验验证。权重预测结果表明，XGBoost 能够较好拟合离线反解得到的样本级融合权重；总体定位性能、CDF 曲线和典型轨迹结果表明，本文方法相比 AoA-only、PDR-only 和 EKF 能够同时降低平均误差和高分位误差，并提升连续轨迹的贴合稳定性；不同信噪比实验进一步说明该方法在观测质量变化时具有较好的鲁棒性；消融实验则验证了样本级权重预测模块对定位精度提升的直接贡献。上述结果表明，结合物理特征、一致性距离和 XGBoost 回归的监督权重学习方法能够有效提升复杂城市环境下的单目标自适应融合定位性能。
